\chapter{ناسیونالیسم و راست افراطی}

\begin{tcolorbox}[
    enhanced,
    colback=lightbg,
    colframe=leftred,
    boxrule=2pt,
    arc=8pt,
    fonttitle=\bfseries\large,
    title={\rl{نقل‌قول آغازین}},
    attach boxed title to top center={yshift=-3mm},
    boxed title style={colback=leftred, colframe=leftred}
]
\begin{center}
\Large\textit{«ملت‌ها ساخته می‌شوند، زاده نمی‌شوند.»}

\vspace{3mm}
\normalsize
--- ارنست گلنر، \textit{ملت‌ها و ناسیونالیسم}، ۱۹۸۳
\end{center}
\end{tcolorbox}

\vspace{5mm}

\section*{درآمد}

\begin{multicols}{2}
ناسیونالیسم یکی از پیچیده‌ترین پدیده‌های سیاسی مدرن است. این جریان فکری در دو قرن گذشته هم امپراتوری‌ها را فروپاشیده، هم کشورها را یکپارچه کرده، هم جنبش‌های آزادی‌بخش الهام‌بخشیده، و هم فاجعه‌بارترین جنایات تاریخ را توجیه کرده است.

آیا ناسیونالیسم جریانی راست است یا چپ؟ در این فصل خواهیم دید که پاسخ ساده نیست. ناسیونالیسم در تاریخ خود هم با لیبرالیسم و جمهوری‌خواهی همراه بوده، هم با سوسیالیسم، و هم با محافظه‌کاری و فاشیسم. اما در دهه‌های اخیر، «راست افراطی» به طور خاص با اشکال معینی از ناسیونالیسم پیوند خورده است که موضوع اصلی این فصل است.
\end{multicols}

%═══════════════════════════════════════════════════════════════
\section{ناسیونالیسم چیست؟}
%═══════════════════════════════════════════════════════════════

\subsection{تعریف و انواع}

\begin{tcolorbox}[
    enhanced,
    colback=rightblue!8,
    colframe=rightblue,
    boxrule=1.5pt,
    arc=5pt,
    fonttitle=\bfseries,
    title={\rl{🔵 تعریف کلیدی: ناسیونالیسم}}
]
\textbf{ناسیونالیسم} عقیده‌ای است که بر این باور استوار است که جهان به \textbf{ملت‌ها}ی متمایز تقسیم شده و هر ملتی حق \textbf{خودمختاری سیاسی} (معمولاً در قالب دولت-ملت) دارد. ناسیونالیسم هم یک \textbf{احساس} است (تعلق به ملت) و هم یک \textbf{ایدئولوژی} (ملت باید دولت داشته باشد).
\end{tcolorbox}

\begin{multicols}{2}
اندیشمندان بین دو نوع ناسیونالیسم تمایز می‌گذارند:

\textbf{ناسیونالیسم مدنی (Civic Nationalism):}
\begin{itemize}
    \item ملت بر اساس شهروندی و قانون تعریف می‌شود
    \item هر کسی می‌تواند عضو ملت شود
    \item نمونه: فرانسه، آمریکا
    \item ارزش‌ها: آزادی، برابری، قانون
\end{itemize}

\textbf{ناسیونالیسم قومی (Ethnic Nationalism):}
\begin{itemize}
    \item ملت بر اساس نژاد، خون، یا زبان تعریف می‌شود
    \item عضویت موروثی است، نه اکتسابی
    \item نمونه: آلمان قرن نوزدهم
    \item تأکید بر تبار و فرهنگ مشترک
\end{itemize}

البته در عمل، این تمایز همیشه روشن نیست. حتی ناسیونالیسم فرانسوی که «مدنی» تلقی می‌شود، در برابر مهاجران مسلمان گاه رنگ قومی می‌گیرد.
\end{multicols}

\subsection{ملت: واقعیت یا ساخته؟}

\begin{multicols}{2}
یکی از مهم‌ترین مناظرات در مطالعات ناسیونالیسم این است: آیا ملت‌ها واقعیت‌های طبیعی و ازلی هستند، یا ساخته‌های مدرن؟

\textbf{ازلی‌گرایان (Primordialists):}
\begin{itemize}
    \item ملت‌ها ریشه در تاریخ کهن دارند
    \item پیوندهای خونی و قومی واقعی هستند
    \item احساس ملی طبیعی و غریزی است
\end{itemize}

\textbf{مدرنیست‌ها (Modernists):}
\begin{itemize}
    \item ملت‌ها ساختهٔ دوران مدرن هستند
    \item صنعتی‌شدن، چاپ، و آموزش همگانی ملت‌ها را ساختند
    \item ناسیونالیسم ایدئولوژی بورژوازی است
\end{itemize}

\textbf{ارنست گلنر} استدلال می‌کرد که ناسیونالیسم محصول نیاز جامعهٔ صنعتی به نیروی کار همگن و باسواد است. \textbf{بندیکت اندرسون} ملت‌ها را «جماعات خیالی» نامید: مردم هم‌ملتان‌شان را نمی‌شناسند، اما تصور می‌کنند با آنها پیوند دارند.
\end{multicols}

\begin{tcolorbox}[
    enhanced,
    colback=centergreen!10,
    colframe=centergreen,
    boxrule=1.5pt,
    arc=5pt,
    fonttitle=\bfseries,
    title={\rl{🟢 نقل‌قول: ارنست رنان}}
]
\textit{«ملت، استفتای روزانه است.»}

فیلسوف فرانسوی ارنست رنان (۱۸۸۲) استدلال می‌کرد که ملت نه بر خون و نه بر زبان استوار است، بلکه بر ارادهٔ مشترک برای زندگی با هم. ملت هر روز با تصمیم شهروندان بازآفریده می‌شود.

\hfill --- ارنست رنان، \textit{ملت چیست؟}، ۱۸۸۲
\end{tcolorbox}

%═══════════════════════════════════════════════════════════════
\section{ریشه‌های تاریخی ناسیونالیسم}
%═══════════════════════════════════════════════════════════════

\subsection{انقلاب فرانسه و زایش ناسیونالیسم}

\begin{multicols}{2}
پیش از انقلاب فرانسه، وفاداری مردم به پادشاه، دین، یا منطقه بود، نه به «ملت». انقلاب فرانسه مفهوم \textbf{حاکمیت ملی} را زایاند: قدرت نه از خدا و پادشاه، بلکه از ملت سرچشمه می‌گیرد.

«مارسی‌یز»، سرود ملی فرانسه، در ۱۷۹۲ سروده شد: «برخیزید فرزندان میهن!» جنگ‌های انقلابی و ناپلئونی ناسیونالیسم را به سراسر اروپا بردند—هم به عنوان الهام و هم به عنوان واکنش.

در آلمان، اشغال ناپلئونی واکنشی ناسیونالیستی برانگیخت. فیلسوفان رمانتیک مانند \textbf{یوهان گوتفرید هردر} (۱۷۴۴-۱۸۰۳) و \textbf{یوهان گوتلیب فیشته} (۱۷۶۲-۱۸۱۴) از «روح ملی» (Volksgeist) و یگانگی ملت آلمان سخن گفتند.

هردر معتقد بود هر ملتی فرهنگ، زبان، و روح منحصر به فرد خود را دارد و نباید تحت سلطهٔ فرهنگ دیگر باشد. این ایده‌ها بنیان ناسیونالیسم قومی را گذاشتند.
\end{multicols}

\subsection{تایم‌لاین: ظهور ناسیونالیسم در اروپا}

\begin{center}
\begin{tikzpicture}[
    scale=0.9,
    every node/.style={font=\small}
]
% Timeline base
\draw[line width=2pt, color=darkgray] (0,0) -- (16,0);

% Period markers
\foreach \x/\year in {0/1789, 4/1815, 8/1848, 12/1871, 16/1914} {
    \draw[line width=1.5pt, color=darkgray] (\x,-0.2) -- (\x,0.2);
    \node[below] at (\x,-0.3) {\year};
}

% Events - alternating above and below
\node[above, align=center, fill=leftred!20, rounded corners, inner sep=3pt] at (0,0.5) {
    \begin{tabular}{c}
    {\footnotesize\rl{انقلاب}}\\
    {\footnotesize\rl{فرانسه}}
    \end{tabular}
};

\node[below, align=center, fill=rightblue!20, rounded corners, inner sep=3pt] at (4,-0.8) {
    \begin{tabular}{c}
    {\footnotesize\rl{کنگرهٔ وین}}\\
    {\footnotesize\rl{نظم ضدانقلابی}}
    \end{tabular}
};

\node[above, align=center, fill=goldaccent!20, rounded corners, inner sep=3pt] at (8,0.5) {
    \begin{tabular}{c}
    {\footnotesize\rl{بهار ملت‌ها}}\\
    {\footnotesize\rl{انقلاب‌های ۱۸۴۸}}
    \end{tabular}
};

\node[below, align=center, fill=centergreen!20, rounded corners, inner sep=3pt] at (10,-0.8) {
    \begin{tabular}{c}
    {\footnotesize\rl{وحدت ایتالیا}}\\
    {\scriptsize 1861}
    \end{tabular}
};

\node[above, align=center, fill=centergreen!20, rounded corners, inner sep=3pt] at (12,0.5) {
    \begin{tabular}{c}
    {\footnotesize\rl{وحدت آلمان}}\\
    {\footnotesize\rl{امپراتوری جدید}}
    \end{tabular}
};

\node[below, align=center, fill=leftred!20, rounded corners, inner sep=3pt] at (16,-0.8) {
    \begin{tabular}{c}
    {\footnotesize\rl{جنگ جهانی}}\\
    {\footnotesize\rl{اول}}
    \end{tabular}
};

% Title
\node[above, font=\bfseries\large] at (8,2) {\rl{تایم‌لاین: قرن ناسیونالیسم (۱۷۸۹-۱۹۱۴)}};

\end{tikzpicture}
\end{center}

\subsection{بهار ملت‌ها ۱۸۴۸}

\begin{multicols}{2}
در ۱۸۴۸، موجی از انقلاب‌ها اروپا را فراگرفت. از فرانسه تا آلمان، ایتالیا، امپراتوری هابسبورگ، و مجارستان، مردم علیه نظام‌های استبدادی شوریدند.

\textbf{خواسته‌های انقلابیون:}
\begin{itemize}
    \item حکومت قانون و پارلمان
    \item آزادی مطبوعات و انجمن
    \item وحدت ملی (در آلمان و ایتالیا)
    \item استقلال ملی (در مجارستان، لهستان، چک)
\end{itemize}

انقلاب‌های ۱۸۴۸ در کوتاه‌مدت شکست خوردند، اما ناسیونالیسم را به ایدئولوژی غالب قرن نوزدهم تبدیل کردند. ظرف دو دهه، ایتالیا و آلمان متحد شدند.
\end{multicols}

%═══════════════════════════════════════════════════════════════
\section{ناسیونالیسم یکپارچه‌ساز}
%═══════════════════════════════════════════════════════════════

\begin{multicols}{2}
در قرن نوزدهم، ناسیونالیسم نیروی یکپارچه‌سازی بود که ملت-دولت‌های مدرن را ساخت.

\textbf{وحدت ایتالیا (Risorgimento):}
\begin{itemize}
    \item \textbf{جوزپه ماتزینی}: ایدئولوگ جمهوری‌خواه ناسیونالیست
    \item \textbf{جوزپه گاریبالدی}: قهرمان نظامی
    \item \textbf{کامیلو کاوور}: سیاستمدار واقع‌گرا
    \item ایتالیا در ۱۸۶۱ تحت سلطنت ساووی متحد شد
\end{itemize}

ماتزینی معتقد بود هر ملتی رسالت الهی دارد و ایتالیای متحد می‌تواند اروپا را به سوی آزادی رهبری کند. او یکی از بنیان‌گذاران «ایتالیای جوان» و «اروپای جوان» بود.

\textbf{وحدت آلمان:}
\begin{itemize}
    \item \textbf{اتو فون بیسمارک}: «آهن و خون»
    \item سه جنگ: دانمارک (۱۸۶۴)، اتریش (۱۸۶۶)، فرانسه (۱۸۷۰)
    \item امپراتوری آلمان در ۱۸۷۱ اعلام شد
\end{itemize}

بیسمارک ناسیونالیسم را نه به عنوان آرمان، بلکه به عنوان ابزار قدرت به کار گرفت. او گفت: «مسائل بزرگ زمانه با سخنرانی و رأی اکثریت حل نمی‌شوند، بلکه با آهن و خون.»
\end{multicols}

%═══════════════════════════════════════════════════════════════
\section{فاشیسم: ناسیونالیسم افراطی}
%═══════════════════════════════════════════════════════════════

\begin{tcolorbox}[
    enhanced,
    colback=rightblue!8,
    colframe=rightblue,
    boxrule=1.5pt,
    arc=5pt,
    fonttitle=\bfseries,
    title={\rl{🔵 تعریف کلیدی: فاشیسم}}
]
\textbf{فاشیسم} جنبش سیاسی افراطی قرن بیستم است که مشخصه‌های زیر را دارد:
\begin{itemize}
    \item \textbf{ناسیونالیسم افراطی}: ملت بالاتر از همه چیز
    \item \textbf{ضد لیبرالیسم}: رد فردگرایی، دموکراسی، حقوق بشر
    \item \textbf{ضد سوسیالیسم}: رد مبارزهٔ طبقاتی و برابری‌خواهی
    \item \textbf{اقتدارگرایی}: رهبر کاریزماتیک، دولت توتالیتر
    \item \textbf{خشونت}: جنگ و استعمار به عنوان فضیلت
    \item \textbf{کورپوراتیسم}: سازماندهی اقتصاد توسط دولت
\end{itemize}
\end{tcolorbox}

\subsection{ایتالیای موسولینی}

\begin{multicols}{2}
\textbf{بنیتو موسولینی} (۱۸۸۳-۱۹۴۵) در جوانی سوسیالیست بود و سردبیر روزنامهٔ حزب سوسیالیست ایتالیا. اما در جنگ جهانی اول، از جنگ حمایت کرد و از حزب اخراج شد.

در ۱۹۱۹ او «فاشی‌های رزم» را تأسیس کرد—گروه‌های شبه‌نظامی که کمونیست‌ها و سوسیالیست‌ها را می‌زدند. در ۱۹۲۱ حزب ملی فاشیست تأسیس شد.

در اکتبر ۱۹۲۲، موسولینی با «راهپیمایی به رم» تهدید کرد. پادشاه، از ترس جنگ داخلی، او را به نخست‌وزیری دعوت کرد. تا ۱۹۲۵، موسولینی دیکتاتور بود.

\textbf{ویژگی‌های فاشیسم ایتالیایی:}
\begin{itemize}
    \item دولت توتالیتر: «همه چیز در دولت، هیچ چیز بیرون دولت»
    \item کورپوراتیسم: سازماندهی اقتصاد بر اساس صنف
    \item امپریالیسم: تهاجم به اتیوپی (۱۹۳۵)
    \item اتحاد با نازیسم
\end{itemize}
\end{multicols}

\subsection{آیا فاشیسم راست است؟}

\begin{tcolorbox}[
    enhanced,
    colback=leftred!8,
    colframe=leftred,
    boxrule=1.5pt,
    arc=5pt,
    fonttitle=\bfseries,
    title={\rl{🔴 موضوع اختلافی: فاشیسم در کجای طیف؟}}
]
\begin{multicols}{2}
\textbf{فاشیسم راست است چون:}
\begin{itemize}
    \item ضد سوسیالیسم و کمونیسم بود
    \item از مالکیت خصوصی حمایت می‌کرد
    \item با نخبگان اقتصادی متحد شد
    \item سلسله‌مراتب طبیعی را می‌پذیرفت
    \item ناسیونالیست و امپریالیست بود
\end{itemize}

\columnbreak

\textbf{فاشیسم راست نیست چون:}
\begin{itemize}
    \item موسولینی سوسیالیست بود
    \item فاشیسم ضد سرمایه‌داری لیبرال بود
    \item دولت بر اقتصاد کنترل داشت
    \item ضد فردگرایی بود
    \item خود را «نه راست، نه چپ» می‌نامید
\end{itemize}
\end{multicols}

\vspace{3mm}
\textbf{نتیجه:} اکثر مورخان فاشیسم را در «راست افراطی» قرار می‌دهند، اما تأکید می‌کنند که این جریان با راست سنتی (محافظه‌کاری، لیبرالیسم کلاسیک) تفاوت‌های بنیادین دارد.
\end{tcolorbox}

%═══════════════════════════════════════════════════════════════
\section{نازیسم}
%═══════════════════════════════════════════════════════════════

\begin{tcolorbox}[
    enhanced,
    colback=leftred!10,
    colframe=leftred!80!black,
    boxrule=2pt,
    arc=5pt,
    fonttitle=\bfseries,
    title={\rl{⚠️ هشدار تاریخی}}
]
این بخش به بررسی نازیسم به عنوان پدیده‌ای تاریخی می‌پردازد. هدف، درک چگونگی ظهور این ایدئولوژی و جلوگیری از تکرار آن است. هیچ بخشی از این تحلیل نباید به عنوان تأیید یا توجیه نازیسم تفسیر شود.
\end{tcolorbox}

\subsection{ایدئولوژی نازیسم}

\begin{multicols}{2}
\textbf{ناسیونال سوسیالیسم} (نازیسم) که توسط \textbf{آدولف هیتلر} (۱۸۸۹-۱۹۴۵) رهبری می‌شد، با فاشیسم ایتالیایی شباهت‌هایی داشت اما تفاوت مهمی هم داشت: \textbf{نژادپرستی بیولوژیک}.

\textbf{اصول نازیسم:}
\begin{itemize}
    \item \textbf{نژاد برتر}: آریایی‌ها (به‌ویژه آلمانی‌ها) نژاد برتر هستند
    \item \textbf{فضای زیستی (Lebensraum)}: آلمان حق دارد سرزمین‌های شرقی را فتح کند
    \item \textbf{یهودستیزی}: یهودیان دشمن نژادی ملت آلمان هستند
    \item \textbf{اصل رهبری (Führerprinzip)}: اطاعت مطلق از هیتلر
    \item \textbf{جامعهٔ نژادی (Volksgemeinschaft)}: وحدت ملی بر اساس خون
\end{itemize}

نازیسم ترکیبی بود از ناسیونالیسم آلمانی، داروینیسم اجتماعی، یهودستیزی قدیمی اروپایی، و نظریات نژادی شبه‌علمی.
\end{multicols}

\subsection{مسیر به قدرت}

\begin{multicols}{2}
هیتلر در ۱۹۲۰ به حزب کارگران آلمان پیوست که بعداً به «حزب ناسیونال سوسیالیست کارگران آلمان» (NSDAP) تغییر نام داد. پس از کودتای نافرجام ۱۹۲۳ (پوتش آبجوخانه)، هیتلر زندانی شد و در زندان کتاب «نبرد من» را نوشت.

در دههٔ ۱۹۳۰، بحران اقتصادی بزرگ و بی‌ثباتی جمهوری وایمار فرصت نازی‌ها را فراهم کرد. در انتخابات ۱۹۳۲، نازی‌ها بزرگ‌ترین حزب رایشتاگ شدند (۳۷٪). در ژانویهٔ ۱۹۳۳، هیتلر صدراعظم شد—نه با کودتا، بلکه با انتصاب قانونی.

ظرف چند ماه، هیتلر دموکراسی را نابود کرد: آتش‌سوزی رایشتاگ بهانه شد، احزاب ممنوع شدند، اتحادیه‌ها منحل شدند، و آلمان به دیکتاتوری تبدیل شد.
\end{multicols}

\subsection{هولوکاست}

\begin{multicols}{2}
نازیسم به فاجعه‌بارترین جنایت تاریخ بشر انجامید: \textbf{هولوکاست} (شوآه)—قتل سیستماتیک حدود شش میلیون یهودی اروپایی، به علاوهٔ میلیون‌ها کولی، معلول، همجنس‌گرا، و دیگر «نامطلوبان».

هولوکاست تنها کشتار نبود؛ \textbf{صنعتی‌سازی مرگ} بود. اردوگاه‌های مرگ مانند آشویتس، تربلینکا، و سوبیبور با کارایی بوروکراتیک کار می‌کردند.

چگونه ممکن شد؟ این پرسشی است که هنوز پاسخ کامل ندارد. \textbf{هانا آرنت} از «ابتذال شر» سخن گفت: جنایتکاران نازی نه هیولا، بلکه بوروکرات‌های معمولی بودند که بدون تفکر، دستور اجرا می‌کردند.
\end{multicols}

%═══════════════════════════════════════════════════════════════
\section{راست افراطی پس از جنگ جهانی دوم}
%═══════════════════════════════════════════════════════════════

\subsection{نئونازیسم و نفی هولوکاست}

\begin{multicols}{2}
پس از ۱۹۴۵، فاشیسم و نازیسم در بیشتر کشورها ممنوع شدند. اما جریان‌های نئونازی و نئوفاشیست در حاشیه زنده ماندند.

\textbf{ویژگی‌های نئونازیسم:}
\begin{itemize}
    \item ستایش هیتلر و رایش سوم
    \item یهودستیزی و نژادپرستی
    \item نفی یا کوچک‌شماری هولوکاست
    \item استفاده از نمادهای نازی
    \item گروه‌های خشونت‌طلب
\end{itemize}

\textbf{نفی هولوکاست} (Holocaust Denial) ادعا می‌کند که هولوکاست یا رخ نداده یا بزرگ‌نمایی شده است. این ادعا توسط تمام مورخان جدی رد شده و در بسیاری از کشورها جرم است.
\end{multicols}

\subsection{راست رادیکال اروپایی}

\begin{multicols}{2}
در دهه‌های اخیر، احزابی که «راست رادیکال» یا «راست پوپولیست» نامیده می‌شوند، در اروپا رشد کرده‌اند. این احزاب با نئونازیسم تفاوت دارند:

\textbf{ویژگی‌های راست رادیکال:}
\begin{itemize}
    \item در چارچوب دموکراسی فعالیت می‌کنند
    \item خشونت را (رسماً) رد می‌کنند
    \item از نژادپرستی آشکار اجتناب می‌کنند
    \item بر «هویت فرهنگی» تأکید می‌کنند
    \item ضد مهاجرت و ضد اسلام هستند
    \item یوروسپتیک (مخالف اتحادیهٔ اروپا) هستند
\end{itemize}

نمونه‌ها: جبههٔ ملی / تجمع ملی (فرانسه)، لیگ (ایتالیا)، آلترناتیو برای آلمان (AfD)، فیدس (مجارستان).
\end{multicols}

\subsection{تمایز مفهومی}

\begin{table}[H]
\centering
\begin{tabular}{|>{\columncolor{rightblue!10}}r|p{3cm}|p{3cm}|p{3cm}|}
\hline
\rowcolor{rightblue!30}
\textbf{معیار} & \textbf{راست افراطی} & \textbf{راست رادیکال} & \textbf{راست میانه} \\
\hline
\textbf{دموکراسی} & ضد دموکراسی & در چارچوب دموکراسی & دموکرات \\
\hline
\textbf{خشونت} & توجیه خشونت & رسماً رد خشونت & رد خشونت \\
\hline
\textbf{نژادپرستی} & آشکار & پنهان/فرهنگی & رد \\
\hline
\textbf{ناسیونالیسم} & نژادی/قومی & فرهنگی/هویتی & مدنی/میانه \\
\hline
\textbf{نمونه} & نئونازی‌ها & لوپن، اوربان & CDU آلمان \\
\hline
\end{tabular}
\caption{طیف راست: از میانه تا افراط}
\end{table}

%═══════════════════════════════════════════════════════════════
\section{موضوعات محوری راست افراطی معاصر}
%═══════════════════════════════════════════════════════════════

\subsection{مهاجرت و اسلام‌هراسی}

\begin{multicols}{2}
موضوع مهاجرت، به‌ویژه مهاجرت از کشورهای مسلمان، موضوع محوری راست افراطی معاصر است.

\textbf{استدلال‌های راست افراطی:}
\begin{itemize}
    \item اسلام با فرهنگ اروپایی ناسازگار است
    \item مهاجران مسلمان «جایگزین بزرگ» جمعیت اروپایی می‌شوند
    \item تروریسم اسلامی تهدید وجودی است
    \item مهاجرت اقتصاد را تضعیف می‌کند
\end{itemize}

\textbf{نظریهٔ «جایگزینی بزرگ»} (Grand Remplacement) ادعا می‌کند که نخبگان عمداً جمعیت سفید اروپایی را با مهاجران غیراروپایی جایگزین می‌کنند. این نظریه توسط نویسندهٔ فرانسوی رنو کامو مطرح شده و در محافل راست افراطی رواج یافته است.

منتقدان این نظریه را «توطئه‌باوری نژادپرستانه» می‌دانند که واقعیت‌های جمعیت‌شناختی را تحریف می‌کند.
\end{multicols}

\subsection{هویت و فرهنگ}

\begin{multicols}{2}
راست افراطی معاصر کمتر از «نژاد» و بیشتر از «فرهنگ» و «هویت» سخ
Claude Opus 4.5
continue


Thought Process
expand
latex
% ادامه فصل ۱۲: ناسیونالیسم و راست افراطی

\subsection{هویت و فرهنگ}

\begin{multicols}{2}
راست افراطی معاصر کمتر از «نژاد» و بیشتر از «فرهنگ» و «هویت» سخن می‌گوید. این تغییر زبان، استراتژیک است: نژادپرستی آشکار غیرقابل قبول است، اما دفاع از «هویت فرهنگی» مشروع‌تر به نظر می‌رسد.

\textbf{نژادپرستی فرهنگی:}
\begin{itemize}
    \item فرهنگ‌ها ذاتاً متفاوت و ناسازگارند
    \item هر فرهنگی باید در سرزمین خود بماند
    \item چندفرهنگی شکست خورده است
    \item اروپا باید مسیحی/سکولار بماند
\end{itemize}

منتقدان می‌گویند این «نژادپرستی بدون نژاد» است: همان محتوای نژادپرستانه با بسته‌بندی جدید. مدافعان می‌گویند دفاع از فرهنگ خودی حق طبیعی هر ملتی است.

\textbf{هویت‌گرایی} (Identitarianism) جنبشی است که از «نسل هویت» (Génération Identitaire) فرانسه سرچشمه گرفته. این جنبش از «نوزایش اروپایی» و حفظ «هویت اروپایی» دفاع می‌کند و علیه «اسلامیزاسیون» مبارزه می‌کند.
\end{multicols}

\subsection{ضد جهانی‌شدن و یوروسپتیسیزم}

\begin{multicols}{2}
راست افراطی معمولاً مخالف جهانی‌شدن و نهادهای فراملی است:

\textbf{ضد جهانی‌شدن:}
\begin{itemize}
    \item جهانی‌شدن به نفع نخبگان بی‌ریشه است
    \item شرکت‌های چندملیتی ملت‌ها را تضعیف می‌کنند
    \item تجارت آزاد، کارگران بومی را بیکار می‌کند
    \item فرهنگ جهانی، فرهنگ‌های ملی را می‌بلعد
\end{itemize}

\textbf{یوروسپتیسیزم:}
\begin{itemize}
    \item اتحادیهٔ اروپا حاکمیت ملی را نقض می‌کند
    \item بروکسل دموکراتیک نیست
    \item یورو به ضرر اقتصادهای ضعیف‌تر است
    \item مرزهای باز، امنیت را تهدید می‌کند
\end{itemize}

\textbf{برگزیت} (خروج بریتانیا از اتحادیهٔ اروپا، ۲۰۱۶) بزرگ‌ترین پیروزی یوروسپتیسیزم بود. شعار «بازپس‌گیری کنترل» (Take Back Control) ترکیبی از ناسیونالیسم، ضد مهاجرت، و ضد بروکسل بود.
\end{multicols}

%═══════════════════════════════════════════════════════════════
\section{شخصیت‌های کلیدی معاصر}
%═══════════════════════════════════════════════════════════════

\begin{tcolorbox}[
    enhanced,
    colback=gray!5,
    colframe=darkgray,
    boxrule=2pt,
    arc=10pt,
    shadow={2mm}{-2mm}{0mm}{black!30},
    fonttitle=\bfseries\large,
    title={\rl{👤 پرترهٔ سیاستمدار: مارین لوپن}}
]
\begin{multicols}{2}
\textbf{مارین لوپن} (متولد ۱۹۶۸)، رهبر حزب «تجمع ملی» (سابقاً جبههٔ ملی) فرانسه.

\textbf{مسیر سیاسی:}
\begin{itemize}
    \item دختر ژان ماری لوپن، بنیان‌گذار جبهه ملی
    \item رهبری حزب از ۲۰۱۱
    \item «زدایش سمیت» (Dédiabolisation) از حزب
    \item نامزد ریاست جمهوری ۲۰۱۷ و ۲۰۲۲
\end{itemize}

\columnbreak

\textbf{مواضع:}
\begin{itemize}
    \item محدودیت شدید مهاجرت
    \item خروج از یورو (سابقاً)
    \item «اولویت ملی» در اشتغال و رفاه
    \item ضد اسلام سیاسی
    \item حمایت‌گرایی اقتصادی
\end{itemize}

لوپن تلاش کرده حزب را از گذشتهٔ یهودستیز و نئوفاشیستی پدرش جدا کند.
\end{multicols}
\end{tcolorbox}

\vspace{5mm}

\begin{tcolorbox}[
    enhanced,
    colback=gray!5,
    colframe=darkgray,
    boxrule=2pt,
    arc=10pt,
    shadow={2mm}{-2mm}{0mm}{black!30},
    fonttitle=\bfseries\large,
    title={\rl{👤 پرترهٔ سیاستمدار: ویکتور اوربان}}
]
\begin{multicols}{2}
\textbf{ویکتور اوربان} (متولد ۱۹۶۳)، نخست‌وزیر مجارستان از ۲۰۱۰.

\textbf{مسیر سیاسی:}
\begin{itemize}
    \item در جوانی لیبرال ضد کمونیست
    \item بنیان‌گذار فیدس (۱۹۸۸)
    \item نخست‌وزیر ۱۹۹۸-۲۰۰۲
    \item چرخش به راست از ۲۰۱۰
\end{itemize}

\columnbreak

\textbf{مواضع:}
\begin{itemize}
    \item «دموکراسی غیرلیبرال»
    \item ضد مهاجرت شدید
    \item ارزش‌های مسیحی و خانواده
    \item ضد بروکسل و سوروس
    \item روابط نزدیک با روسیه و چین
\end{itemize}

اوربان الگویی برای راست پوپولیست جهان شده است.
\end{multicols}
\end{tcolorbox}

%═══════════════════════════════════════════════════════════════
\section{نقد و تحلیل}
%═══════════════════════════════════════════════════════════════

\subsection{چرا راست افراطی رشد می‌کند؟}

\begin{multicols}{2}
در دهه‌های اخیر، احزاب راست افراطی و رادیکال در سراسر جهان رشد کرده‌اند. تبیین‌های مختلفی برای این پدیده ارائه شده:

\textbf{تبیین اقتصادی:}
\begin{itemize}
    \item جهانی‌شدن، طبقهٔ کارگر سفید را متضرر کرده
    \item صنعت‌زدایی و بیکاری
    \item نابرابری فزاینده
    \item رقابت با مهاجران بر سر شغل و مسکن
\end{itemize}

\textbf{تبیین فرهنگی:}
\begin{itemize}
    \item واکنش به تغییرات فرهنگی سریع
    \item احساس از دست دادن هویت
    \item «محرومیت نسبی» اکثریت قدیمی
    \item مقاومت در برابر چندفرهنگی و «صحت سیاسی»
\end{itemize}

\textbf{تبیین سیاسی:}
\begin{itemize}
    \item بحران اعتماد به نخبگان
    \item احساس نمایندگی نشدن
    \item شکست احزاب سنتی
    \item قطبی‌شدن رسانه‌ای
\end{itemize}

احتمالاً ترکیبی از همهٔ این عوامل در کار است. پژوهش‌ها نشان می‌دهد که رأی‌دهندگان راست افراطی هم «بازندگان جهانی‌شدن» هستند و هم کسانی که از تغییر فرهنگی ناراضی‌اند.
\end{multicols}

\subsection{محرومیت یا نژادپرستی؟}

\begin{tcolorbox}[
    enhanced,
    colback=leftred!8,
    colframe=leftred,
    boxrule=1.5pt,
    arc=5pt,
    fonttitle=\bfseries,
    title={\rl{🔴 مناظره: رأی‌دهندگان راست افراطی}}
]
\begin{multicols}{2}
\textbf{دیدگاه اول: قربانیان نظام}
\begin{itemize}
    \item این مردم توسط نخبگان رها شده‌اند
    \item نگرانی‌های اقتصادی‌شان مشروع است
    \item به آنها باید گوش داد، نه تحقیرشان کرد
    \item چپ باید طبقهٔ کارگر را بازپس گیرد
\end{itemize}

\columnbreak

\textbf{دیدگاه دوم: نژادپرستی توده‌ای}
\begin{itemize}
    \item مسئله اقتصادی نیست، نژادپرستی است
    \item نباید نژادپرستی را عادی‌سازی کرد
    \item «درک» آنها یعنی پذیرش نژادپرستی
    \item باید با این جریان مبارزه کرد
\end{itemize}
\end{multicols}

\vspace{3mm}
\textbf{موضع میانه:} احتمالاً هر دو عامل دخیل‌اند. محرومیت اقتصادی زمینهٔ پذیرش پیام‌های نژادپرستانه را فراهم می‌کند، اما خود نژادپرستی هم عامل مستقلی است.
\end{tcolorbox}

%═══════════════════════════════════════════════════════════════
\section*{خلاصهٔ فصل}
%═══════════════════════════════════════════════════════════════

\begin{tcolorbox}[
    enhanced,
    colback=lightbg,
    colframe=darkgray,
    boxrule=2pt,
    arc=8pt,
    fonttitle=\bfseries\large,
    title={\rl{📝 خلاصهٔ فصل دوازدهم}}
]
\begin{itemize}
    \item \textbf{ناسیونالیسم} عقیده‌ای است که جهان را به ملت‌های متمایز تقسیم می‌کند و برای هر ملت حق خودمختاری قائل است. ناسیونالیسم می‌تواند \textbf{مدنی} (شهروندی) یا \textbf{قومی} (خونی) باشد.
    
    \item ناسیونالیسم در \textbf{انقلاب فرانسه} زاده شد و در قرن نوزدهم به ایدئولوژی غالب تبدیل شد. رمانتیک‌های آلمانی مفهوم «روح ملی» را توسعه دادند.
    
    \item \textbf{فاشیسم} (ایتالیا) و \textbf{نازیسم} (آلمان) اشکال افراطی ناسیونالیسم در قرن بیستم بودند. نازیسم با افزودن \textbf{نژادپرستی بیولوژیک} به فاشیسم، به \textbf{هولوکاست} انجامید.
    
    \item پس از جنگ جهانی دوم، \textbf{نئونازیسم} در حاشیه ماند، اما \textbf{راست رادیکال} جدید در چارچوب دموکراسی رشد کرد.
    
    \item موضوعات محوری راست افراطی معاصر: \textbf{ضد مهاجرت}، \textbf{اسلام‌هراسی}، \textbf{هویت فرهنگی}، \textbf{ضد جهانی‌شدن}، و \textbf{یوروسپتیسیزم}.
    
    \item رشد راست افراطی ترکیبی از عوامل \textbf{اقتصادی} (محرومیت)، \textbf{فرهنگی} (واکنش به تغییر)، و \textbf{سیاسی} (بی‌اعتمادی به نخبگان) است.
\end{itemize}
\end{tcolorbox}

%═══════════════════════════════════════════════════════════════
\section*{پرسش‌های تأملی}
%═══════════════════════════════════════════════════════════════

\begin{tcolorbox}[
    enhanced,
    colback=centergreen!8,
    colframe=centergreen,
    boxrule=1.5pt,
    arc=5pt,
    fonttitle=\bfseries,
    title={\rl{❓ پرسش‌هایی برای تأمل و بحث}}
]
\begin{enumerate}
    \item آیا ناسیونالیسم ذاتاً خطرناک است، یا می‌تواند شکل‌های مثبت (مانند جنبش‌های آزادی‌بخش) داشته باشد؟
    
    \item تفاوت «دفاع از فرهنگ» و «نژادپرستی» چیست؟ آیا می‌توان از هویت فرهنگی دفاع کرد بدون اینکه به نژادپرستی افتاد؟
    
    \item چرا فاشیسم و نازیسم دقیقاً در بین دو جنگ جهانی ظهور کردند؟ چه شرایطی این جریان‌ها را ممکن کرد؟
    
    \item آیا اتحادیهٔ اروپا پروژه‌ای ضد ناسیونالیستی است؟ آیا می‌توان هم اروپایی بود و هم ناسیونالیست؟
    
    \item چگونه باید با رشد راست افراطی مقابله کرد؟ سرکوب، گفتگو، یا پاسخ به نگرانی‌های مشروع؟
    
    \item آیا «دموکراسی غیرلیبرال» اوربان واقعاً دموکراسی است، یا شکلی از اقتدارگرایی؟
\end{enumerate}
\end{tcolorbox}

\vspace{10mm}
\begin{center}
\rule{0.6\textwidth}{1pt}

\vspace{3mm}
\textit{«ناسیونالیسم بیماری کودکی بشریت است. سرخک نوع بشر.»}

\vspace{2mm}
--- آلبرت اینشتین

\vspace{3mm}
\rule{0.6\textwidth}{1pt}
\end{center}

